\documentclass[border=10pt]{standalone}
\usepackage{fontspec}
\setmainfont{Noto Serif}[Ligatures=TeXOff]
\XeTeXgenerateactualtext=1
\newfontfamily\bengalifont{Noto Serif Bengali}[Script=Bengali, NFSSFamily=bnfam]
\usepackage[Latin,Bengali]{ucharclasses}
% group-free font switching (as in header-common.tex): robust inside braces, cells and quotes
\makeatletter
\def\BN@saved{\rmdefault}
\def\BN@bn{bnfam}
\newcommand\BNin{\edef\BN@tmp{\f@family}\ifx\BN@tmp\BN@bn\else\let\BN@saved\BN@tmp\fi\fontfamily{bnfam}\selectfont}
\newcommand\BNout{\fontfamily{\BN@saved}\selectfont}
\makeatother
\setTransitionsFor{Bengali}{\BNin}{\BNout}
% Typographic quotes and dashes are not in a ucharclasses class (class 0), so a Bangla word right
% after “ or ‘ would not trigger the Bangla font. Treat them as Latin so that every switch fires.
\XeTeXcharclass"2018=\XeTeXcharclass`A
\XeTeXcharclass"2019=\XeTeXcharclass`A
\XeTeXcharclass"201C=\XeTeXcharclass`A
\XeTeXcharclass"201D=\XeTeXcharclass`A
\XeTeXcharclass"2013=\XeTeXcharclass`A
\XeTeXcharclass"2014=\XeTeXcharclass`A
\XeTeXcharclass"2026=\XeTeXcharclass`A
\usepackage{newunicodechar}
\newunicodechar{।}{{\bengalifont\char"0964}}
\newunicodechar{॥}{{\bengalifont\char"0965}}

\usepackage[table]{xcolor}
\usepackage{booktabs,tabularx,array}
\definecolor{bapilrule}{HTML}{1F2937}
\definecolor{bapilstripe}{HTML}{F3F4F6}
\renewcommand{\arraystretch}{1.28}
\setlength{\tabcolsep}{5pt}

\begin{document}
\begin{minipage}{17cm}
{\large\bfseries সারণি ১৬।}\ {\large নির্দেশক মূল জনবল পরিকল্পনা (পূর্ণকালীন সমতুল্য)}\par\vspace{7pt}
\small
\rowcolors{2}{bapilstripe}{white}
\begin{tabularx}{\linewidth}{@{}>{\hsize=2.000\hsize\raggedright\arraybackslash}X>{\hsize=0.667\hsize\raggedright\arraybackslash}X>{\hsize=0.667\hsize\raggedright\arraybackslash}X>{\hsize=0.667\hsize\raggedright\arraybackslash}X@{}}
\arrayrulecolor{bapilrule}\toprule
\rowcolor{white}\textbf{ভূমিকা} & \textbf{মাস ১-১২} & \textbf{মাস ১৩-২৪} & \textbf{মাস ২৫-৩৬} \\
\midrule
পরিচালক ও প্রধান গবেষক & ১ & ১ & ১ \\
নীতি-গবেষণা প্রধান ও বিশ্লেষক & ১ & ২ & ৩ \\
মেশিন লার্নিং ও বাংলা NLP প্রকৌশলী & ২ & ৩ & ৪ \\
ডেটা প্রকৌশলী & ১ & ২ & ২ \\
সফটওয়্যার ও প্ল্যাটফর্ম (DevSecOps) প্রকৌশলী & ১ & ২ & ২ \\
নিরাপত্তা প্রকৌশলী & ০.৫ & ১ & ১ \\
শাসন, আইন ও নৈতিকতা বিশ্লেষক & ১ & ১ & ২ \\
ব্যবহারকারী অভিজ্ঞতা ও প্রবেশযোগ্যতা নকশাকার & ০.৫ & ১ & ১ \\
কর্মসূচি ও অংশীদারিত্ব সমন্বয়ক & ১ & ১ & ১ \\
প্রশিক্ষণ ও প্রচার কর্মকর্তা & ০ & ১ & ২ \\
\textbf{মোট মূল জনবল} & \textbf{৯} & \textbf{১৫} & \textbf{১৯} \\
\bottomrule
\end{tabularx}
\par\vspace{5pt}{\footnotesize\itshape খণ্ডকালীন ERRC সদস্য, উপদেষ্টা, অতিথি ফেলো ও স্নাতকোত্তর শিক্ষার্থী অন্তর্ভুক্ত নয়। নির্দেশক ব্যয় সারণি ৪৯-এ; চূড়ান্ত ব্যয় সম্ভাব্যতা সমীক্ষায় (D1.3) নির্ধারিত হবে।}
\par\vspace{3pt}{\scriptsize\color{gray} উৎস: BAPIL টেকনিক্যাল ব্লুপ্রিন্ট, সংস্করণ ১.০ (২০২৬)। CC BY 4.0।}
\end{minipage}
\end{document}
